\chapter{Pendahuluan}
\section{Latar Belakang}
Jelaskan konteks, fakta pendukung, kesenjangan penelitian, dan alasan masalah
perlu diteliti. Contoh sitasi naratif dapat ditulis dengan
\textcite{knuth1984texbook}, sedangkan sitasi dalam kurung menggunakan
\parencite{lamport1994latex}.

\section{Rumusan Masalah}
Tuliskan pertanyaan penelitian secara terukur dan selaras dengan tujuan.

\section{Tujuan}
Tuliskan hasil ilmiah yang hendak dicapai oleh penelitian.

\section{Manfaat}
Jelaskan manfaat akademik dan praktis tanpa menjanjikan hasil yang belum diuji.

\section{Batasan Masalah}
Nyatakan objek, data, metode, serta kondisi yang berada di luar lingkup.

\section{Sistematika Pembahasan}
Proposal terdiri atas Pendahuluan, Landasan Kepustakaan, dan Metodologi Penelitian.
